\subsection{SUBALGEBRAS REDUCTIVE IN A LIE ALGEBRA}

DEFINITION 5. Let g be a Lie algebra and h a Lie subalgebra of g. h is called reductive in g if the representation \(x \mapsto ad_{g} x\) of h is semi-simple.

This representation admits as subrepresentation the adjoint representation of \(\mathfrak{h}\) . Hence, if \(\mathfrak{h}\) is reductive in \(\mathfrak{g},\mathfrak{h}\) is reductive. On the other hand, to say that a Lie algebra is reductive in itself is equivalent to saying that it is reductive.

PROPOSITION 7. Let g be a Lie algebra, h a subalgebra reductive in g, ρ a representation of g on a vector space V and W the sum of the finite-dimensional subspaces of V which are simple h-modules. Then W is stable under ρ(g).

Let \(W_0\) be a finite-dimensional simple sub- \(\mathfrak{h}\) -module of V. We need to prove that \(\rho(x)(W_0) \subset W\) for all \(x \in \mathfrak{g}\) . Let M denote the vector space \(\mathfrak{g}\) considered as an \(\mathfrak{h}\) -module by means of the representation \(x \mapsto \operatorname{ad}_{\mathfrak{g}} x\) of \(\mathfrak{h}\) on \(\mathfrak{g}\) . Then \(M \otimes_{K} W_0\) is a semi-simple \(\mathfrak{h}\) -module (Corollary 1 to Theorem 4). Let \(\theta\) be the K-linear mapping of \(\mathbf{M} \otimes_{\mathbb{K}} \mathbf{W}_0\) into \(\mathbf{V}\) defined by \(\theta(x \otimes w) = \rho(x)w\) . This is an \(\mathfrak{h}\) -module homomorphism, for if \(y \in \mathfrak{h}\) then:

\[
\begin{array}{r l} \theta ([ y, x ] \otimes w + x \otimes \rho (y) w) & = \rho ([ y, x ]) w + \rho (x) \rho (y) w \\ & = \rho (y) \rho (x) w = \rho (y) \theta (x \otimes w). \end{array}
\]

Hence \(\theta (\mathbf{M}\otimes_{\mathbb{K}}\mathbf{W}_0)\) is a finite-dimensional semi-simple \(\mathfrak{h}\) -module. Hence \(\theta (\mathbf{M}\otimes_{\mathbb{K}}\mathbf{W}_0)\subset \mathbf{W}\) , that is \(\rho (x)(\mathbf{W}_0)\subset \mathbf{W}\) for all \(x\in \mathfrak{g}\) .

COROLLARY 1. Let g be a Lie algebra, h a subalgebra reductive in g and \(\rho\) a finite-dimensional semi-simple representation of g. Then the restriction of \(\rho\) to h is semi-simple.

It suffices to consider the case where \(\rho\) is simple. We adopt the notation V, W of Proposition 4. Let \(W_{1}\) be a subspace of V minimal among the nonzero subspaces stable under \(\rho(\mathfrak{h})\) . Then \(W_{1} \subset W\) , hence \(W \neq \{0\}\) and hence \(W = V\) .

COROLLARY 2. Let g be a Lie algebra, h a subalgebra reductive in g and k a subalgebra of h reductive in h. Then k is reductive in g.

The representation \(x \mapsto \operatorname{ad}_{\mathfrak{g}} x\) of \(\mathfrak{h}\) on \(\mathfrak{g}\) is semi-simple and hence its restriction to \(k\) is semi-simple (Corollary 1).

